\documentclass[a4paper,12pt]{article}

\usepackage[T2A]{fontenc}
\usepackage[utf8]{inputenc}
\usepackage[russian]{babel}

\usepackage{amsmath,amssymb,amsthm,mathtools}
\usepackage{helvet}
\renewcommand{\familydefault}{\sfdefault}
\usepackage{icomma}
\usepackage[margin=2.1cm]{geometry}
\usepackage{microtype}
\usepackage{tikz}
\usepackage{tkz-euclide}
\usepackage{pgfplots}
\pgfplotsset{compat=1.18}
\usepackage{tabularx,booktabs,longtable}
\usepackage{enumitem}
\usepackage{hyperref}
\usepackage{parskip}
\usepackage{fancyhdr}
\usepackage{setspace}
\usepackage{xcolor}

\definecolor{accent}{HTML}{008080}
\definecolor{darkaccent}{HTML}{004D4D}
\definecolor{textgray}{HTML}{333333}

\hypersetup{
  colorlinks=true,
  linkcolor=darkaccent,
  urlcolor=darkaccent,
  citecolor=darkaccent
}

\onehalfspacing
\setlength{\parindent}{0pt}
\setlength{\headheight}{15pt}
\pagestyle{fancy}
\fancyhf{}
\fancyhead[L]{\textcolor{darkaccent}{Чек-лист по теме занятия}}
\fancyhead[R]{\textcolor{darkaccent}{8 класс}}
\fancyfoot[C]{\thepage}

\renewcommand{\arraystretch}{1.35}
\setlist[itemize]{leftmargin=1.8em}
\setlist[enumerate]{leftmargin=2em}

\begin{document}
\color{textgray}

% ---------------- TITUL ----------------
\begin{titlepage}
\centering
\vspace*{2.5cm}

{\Huge\bfseries \textcolor{darkaccent}{Чек-лист}\par}
\vspace{0.6cm}
{\LARGE\bfseries Графики элементарных функций\\[0.2em]и площади в треугольнике\par}

\vspace{1.8cm}

{\Large Лёвин Артём Александрович\\[0.2em]
эксклюзивно для Софии Кальней\par}

\vspace{0.8cm}
{\large \today\par}

\vfill

{\large\textcolor{accent}{Ключевые темы занятия: прямая, парабола, гипербола,\\
знаки коэффициентов, вершина параболы, подобие и площади.}\par}

\vfill

\begin{tikzpicture}[remember picture,overlay]
\draw[line width=1.2pt, color=accent]
  ([xshift=1.8cm,yshift=3.3cm]current page.south west)
    .. controls ([xshift=5.5cm,yshift=1.8cm]current page.south west)
    and ([xshift=-7.2cm,yshift=4.8cm]current page.south east)
  .. ([xshift=-2.0cm,yshift=2.2cm]current page.south east);

\draw[line width=0.8pt, color=darkaccent]
  ([xshift=2.7cm,yshift=2.35cm]current page.south west)
    .. controls ([xshift=7.0cm,yshift=3.8cm]current page.south west)
    and ([xshift=-5.8cm,yshift=1.0cm]current page.south east)
  .. ([xshift=-2.8cm,yshift=3.1cm]current page.south east);
\end{tikzpicture}
\end{titlepage}

% ---------------- MAIN ----------------
\section*{1. Что нужно уверенно помнить}

\subsection*{1.1. Прямая}
Функция прямой:
\[
y = kx + b.
\]

\begin{itemize}
  \item \(k\) --- угловой коэффициент, он отвечает за наклон.
  \item \(b\) --- ордината точки пересечения с осью \(Oy\), то есть значение \(y\) при \(x=0\).
\end{itemize}

\textbf{Как быстро читать график прямой:}
\begin{itemize}
  \item если \(k>0\), прямая возрастает;
  \item если \(k<0\), прямая убывает;
  \item если прямая пересекает ось \(Oy\) выше нуля, то \(b>0\);
  \item если ниже нуля, то \(b<0\).
\end{itemize}

\textbf{Быстрый способ найти \(k\):}
выберите на графике две удобные точки и посчитайте
\[
k=\frac{\Delta y}{\Delta x}.
\]
Это фактически тангенс угла наклона прямой к положительному направлению оси \(Ox\).

\textbf{Мини-пример:}
если прямая проходит через точки \((0;-2)\) и \((1;1)\), то
\[
k=\frac{1-(-2)}{1-0}=3,\qquad b=-2,
\]
значит, её формула:
\[
y=3x-2.
\]

\subsection*{1.2. Парабола}
Квадратичная функция:
\[
y=ax^2+bx+c,\qquad a\neq 0.
\]

\textbf{Главные параметры:}
\begin{itemize}
  \item \(a\) отвечает за направление ветвей:
  \[
  a>0 \Rightarrow \text{ветви вверх}, \qquad a<0 \Rightarrow \text{ветви вниз}.
  \]
  \item \(c\) --- значение функции при \(x=0\), то есть точка пересечения с осью \(Oy\).
  \item абсцисса вершины:
  \[
  x_{\text{в}}=-\frac{b}{2a};
  \]
  \item ордината вершины:
  \[
  y_{\text{в}}=\frac{4ac-b^2}{4a}.
  \]
\end{itemize}

\textbf{Очень важная формула для знака коэффициента \(b\):}
\[
b=-2ax_{\text{в}}.
\]
Значит, знак \(b\) можно определить по знакам \(a\) и \(x_{\text{в}}\).

\textbf{Как читать знаки коэффициентов по графику параболы:}
\begin{itemize}
  \item \(a\) --- по направлению ветвей;
  \item \(c\) --- по точке пересечения с осью \(Oy\);
  \item \(b\) --- по знаку числа \(-2ax_{\text{в}}\).
\end{itemize}

\textbf{Пример:}
если ветви направлены вверх, вершина расположена правее оси \(Oy\), а график пересекает ось \(Oy\) выше нуля, то
\[
a>0,\qquad x_{\text{в}}>0,\qquad c>0,
\]
следовательно,
\[
b=-2ax_{\text{в}}<0.
\]

\subsection*{1.3. Гипербола вида \(\boldsymbol{y=\frac{k}{x}}\) и \(\boldsymbol{y=\frac{1}{kx}}\)}
\textbf{Базовая форма:}
\[
y=\frac{k}{x},\qquad x\neq 0.
\]

\begin{itemize}
  \item если \(k>0\), ветви расположены в I и III четвертях;
  \item если \(k<0\), ветви расположены во II и IV четвертях.
\end{itemize}

\textbf{Роль коэффициента:}
\begin{itemize}
  \item в записи \(\displaystyle y=\frac{k}{x}\): чем больше \(|k|\), тем дальше ветви от осей;
  \item в записи \(\displaystyle y=\frac{1}{kx}\): чем больше \(|k|\), тем ближе ветви к осям.
\end{itemize}

\textbf{Мини-примеры:}
\[
y=\frac{2}{x}\quad \text{--- ветви дальше от осей, чем у } y=\frac{1}{x};
\]
\[
y=\frac{1}{2x}\quad \text{--- ветви ближе к осям, чем у } y=\frac{1}{x}.
\]

\subsection*{1.4. Что часто спрашивают в заданиях на графики}
\begin{itemize}
  \item установить соответствие между формулой и графиком;
  \item определить знак коэффициентов;
  \item найти вершину параболы;
  \item понять, возрастает или убывает прямая;
  \item определить, в каких четвертях лежат ветви гиперболы;
  \item сравнить график с ``эталонным'' графиком и понять, удалён ли он от осей или прижат к ним.
\end{itemize}

% ---------------- VISUALS ----------------
\section*{2. Визуальные опоры}

\subsection*{2.1. Парабола и её ключевые элементы}
\begin{center}
\begin{tikzpicture}
\begin{axis}[
  width=0.78\linewidth,
  height=6.6cm,
  axis lines=middle,
  xmin=-4, xmax=4,
  ymin=-3, ymax=6,
  unit vector ratio=1 1,
  samples=220,
  clip=false,
  xlabel={$x$},
  ylabel={$y$},
  xtick={-3,-2,-1,0,1,2,3},
  ytick={-2,0,2,4},
  ticklabel style={font=\small},
  label style={font=\small}
]
\addplot[domain=-3.2:3.2, line width=1.15pt, color=accent] {x^2-2*x+1};
\addplot[only marks, mark=*, mark size=1.9pt, color=darkaccent] coordinates {(1,0) (0,1)};
\node[font=\small, anchor=south west, text=darkaccent] at (axis cs:1,0) {вершина};
\node[font=\small, anchor=south east, text=darkaccent] at (axis cs:0,1) {$c$};
\end{axis}
\end{tikzpicture}
\end{center}

\textbf{Запомните:}
\begin{itemize}
  \item вершина справа от оси \(Oy\) \(\Rightarrow x_{\text{в}}>0\);
  \item вершина слева от оси \(Oy\) \(\Rightarrow x_{\text{в}}<0\);
  \item по знаку \(x_{\text{в}}\) и знаку \(a\) можно найти знак \(b\).
\end{itemize}

\subsection*{2.2. Схема к задаче на площади}
Ниже приведена аккуратная схема к ключевой геометрической задаче занятия.

\begin{center}
\begin{tikzpicture}[scale=1.05]
\tkzSetUpLine[line width=0.9pt, color=accent]
\tkzSetUpPoint[color=accent, fill=accent, size=3]
\tkzSetUpLabel[color=black, font=\small]

\tkzDefPoint(0,0){A}
\tkzDefPoint(6,0){C}
\tkzDefPoint(1.3,4.1){B}
\tkzDefMidPoint(A,C)\tkzGetPoint{M}
\tkzDefMidPoint(B,M)\tkzGetPoint{K}

\tkzDefLine[parallel=through K](A,B)\tkzGetPoint{kAB}
\tkzInterLL(K,kAB)(A,C)\tkzGetPoint{J}
\tkzInterLL(K,kAB)(B,C)\tkzGetPoint{P}

\tkzDrawSegments(A,B B,C A,C B,M J,P)
\tkzDrawPoints(A,B,C,M,K,J,P)
\tkzLabelPoints[below](A,C,M,J)
\tkzLabelPoints[above](B,K,P)

\draw[dashed, color=darkaccent] (B)--(M);

\node[font=\small, text=darkaccent, anchor=west] at (3.8,3.15) {$JP\parallel AB$};
\node[font=\small, text=darkaccent, anchor=west] at (2.4,1.8) {$K$ --- середина \(BM\)};
\node[font=\small, text=darkaccent, anchor=west] at (2.4,1.25) {$M$ --- середина \(AC\)};
\end{tikzpicture}
\end{center}

% ---------------- GEOMETRY ----------------
\section*{3. Ключевая геометрическая задача занятия}

\textbf{Условие.}
В треугольнике \(ABC\) точка \(M\) --- середина \(AC\), поэтому \(BM\) --- медиана.
Точка \(K\) --- середина \(BM\).
Через \(K\) проведена прямая, параллельная \(AB\), которая пересекает \(AC\) в точке \(J\), а \(BC\) --- в точке \(P\).
Найти отношение
\[
S_{ABK}:S_{KPCM}.
\]

\subsection*{3.1. Идеи, без которых задача не решается}
\begin{enumerate}
  \item Медиана делит треугольник на два равновеликих треугольника:
  \[
  S_{ABM}=S_{MBC}=\frac12 S_{ABC}.
  \]
  \item Так как \(K\) --- середина \(BM\), то
  \[
  S_{ABK}=\frac12 S_{ABM}=\frac14 S_{ABC}.
  \]
  \item В треугольнике \(ABM\) прямая \(KJ\parallel AB\), а \(K\) --- середина \(BM\), значит, \(J\) --- середина \(AM\).
  \item Так как \(M\) --- середина \(AC\), то
  \[
  AM=\frac12 AC,\qquad AJ=\frac12 AM=\frac14 AC.
  \]
  \item Из подобия треугольников \(CJP\) и \(CAB\):
  \[
  \frac{CP}{CB}=\frac{CJ}{CA}=1-\frac{AJ}{AC}=1-\frac14=\frac34.
  \]
  Значит,
  \[
  \frac{BP}{BC}=\frac14.
  \]
\end{enumerate}

\subsection*{3.2. Удобное вычисление площадей}
Пусть
\[
S_{ABC}=x.
\]
Тогда
\[
S_{ABK}=\frac{x}{4}.
\]

Кроме того,
\[
S_{CJP}=\left(\frac34\right)^2 x=\frac{9x}{16},
\]
потому что площади подобных треугольников относятся как квадраты коэффициента подобия.

В треугольнике \(AKM\) точка \(J\) --- середина \(AM\), следовательно, \(KJ\) --- медиана, поэтому
\[
S_{KJM}=\frac12 S_{AKM}.
\]
Но
\[
S_{AKM}=\frac12 S_{ABM}=\frac12\cdot \frac{x}{2}=\frac{x}{4},
\]
значит,
\[
S_{KJM}=\frac{x}{8}.
\]

Тогда площадь четырёхугольника \(KPCM\) равна
\[
S_{KPCM}=S_{CJP}-S_{KJM}
=\frac{9x}{16}-\frac{x}{8}
=\frac{9x}{16}-\frac{2x}{16}
=\frac{7x}{16}.
\]

Следовательно,
\[
S_{ABK}:S_{KPCM}=\frac{x}{4}:\frac{7x}{16}=4:7.
\]

\subsection*{3.3. Что здесь нужно было заметить}
\begin{itemize}
  \item середины отрезков дают медианы и средние линии;
  \item параллельность почти всегда ведёт к подобию;
  \item площади удобно выражать через одну общую величину;
  \item если видите середину стороны в треугольнике, сразу проверяйте, не появилась ли средняя линия.
\end{itemize}

% ---------------- ALGORITHMS ----------------
\section*{4. Алгоритмы решения}

\subsection*{4.1. Если перед Вами график прямой}
\begin{enumerate}
  \item Определите, возрастает она или убывает.
  \item Найдите точку пересечения с осью \(Oy\) --- это \(b\).
  \item Возьмите две удобные точки и найдите
  \[
  k=\frac{\Delta y}{\Delta x}.
  \]
  \item Запишите формулу \(y=kx+b\).
\end{enumerate}

\subsection*{4.2. Если перед Вами график параболы}
\begin{enumerate}
  \item По направлению ветвей определите знак \(a\).
  \item По пересечению с осью \(Oy\) определите знак \(c\).
  \item Посмотрите, где находится вершина: слева или справа от оси \(Oy\).
  \item Используйте
  \[
  b=-2ax_{\text{в}}
  \]
  и определите знак \(b\).
  \item При необходимости вычислите координаты вершины формулами
  \[
  x_{\text{в}}=-\frac{b}{2a}, \qquad
  y_{\text{в}}=\frac{4ac-b^2}{4a}.
  \]
\end{enumerate}

\subsection*{4.3. Если перед Вами график гиперболы}
\begin{enumerate}
  \item Посмотрите, в каких четвертях лежат ветви:
  \[
  \text{I и III} \Rightarrow k>0,\qquad
  \text{II и IV} \Rightarrow k<0.
  \]
  \item Сравните удалённость ветвей от осей с графиком \(y=\frac1x\).
  \item Определите, какая запись подходит:
  \[
  y=\frac{k}{x}\quad \text{или}\quad y=\frac{1}{kx}.
  \]
\end{enumerate}

\subsection*{4.4. Если задача на площади в треугольнике}
\begin{enumerate}
  \item Проверьте, есть ли медианы, середины сторон, средние линии.
  \item Ищите пары подобных треугольников.
  \item Обозначьте всю площадь или удобную часть через \(x\).
  \item Выражайте остальные площади через \(x\).
  \item В конце составьте отношение и сократите.
\end{enumerate}

% ---------------- MISTAKES ----------------
\section*{5. Типичные ошибки}

\begin{itemize}
  \item Путать коэффициенты \(b\) в прямой \(y=kx+b\) и в параболе \(y=ax^2+bx+c\).
  \item Забывать, что \(c\) в параболе --- это именно пересечение с осью \(Oy\).
  \item Неверно определять знак \(b\), не используя формулу
  \[
  b=-2ax_{\text{в}}.
  \]
  \item Путать графики \(\dfrac{k}{x}\) и \(\dfrac{1}{kx}\).
  \item В задачах на гиперболу смотреть только на знак, но не учитывать удалённость ветвей от осей.
  \item В задачах на площади не использовать равенство площадей при медиане.
  \item При подобии треугольников забывать, что площади относятся как \emph{квадраты} коэффициента подобия.
  \item В геометрии не дорисовывать полезные вспомогательные линии до появления очевидного подобия.
\end{itemize}

% ---------------- SUMMARY TABLE ----------------
\section*{6. Краткая памятка}

\begin{table}[h!]
\centering
\caption{Что за что отвечает}
\begin{tabularx}{\linewidth}{@{}lX@{}}
\toprule
\textbf{Объект} & \textbf{Что важно помнить} \\
\midrule
Прямая \(y=kx+b\) & \(k\) --- наклон, \(b\) --- пересечение с осью \(Oy\). \\
Парабола \(y=ax^2+bx+c\) & \(a\) --- вверх/вниз; \(c\) --- пересечение с \(Oy\); \(x_{\text{в}}=-\dfrac{b}{2a}\). \\
Вершина параболы & \(y_{\text{в}}=\dfrac{4ac-b^2}{4a}\), \quad \(b=-2ax_{\text{в}}\). \\
Гипербола \(y=\dfrac{k}{x}\) & Знак \(k\) задаёт четверти; модуль \(k\) влияет на удалённость от осей. \\
Площади в треугольнике & Медиана делит треугольник на равновеликие части. \\
Подобные треугольники & Стороны относятся как \(m:n\), площади --- как \(m^2:n^2\). \\
\bottomrule
\end{tabularx}
\end{table}

\newpage

% ---------------- TRAINING ----------------
\section*{Тренировочный блок}

\textbf{Задания 1--3 --- базовый уровень.}

\begin{enumerate}
  \item Для прямой, проходящей через точки \((0;3)\) и \((2;-1)\), найдите коэффициенты \(k\) и \(b\), затем запишите её формулу.

  \item По графику параболы известно: ветви направлены вверх, вершина расположена правее оси \(Oy\), график пересекает ось \(Oy\) ниже нуля. Определите знаки коэффициентов \(a\), \(b\), \(c\).

  \item Для функции
  \[
  y=-2x^2+4x+1
  \]
  найдите координаты вершины.

  \item Установите соответствие между формулами и типами графиков:
  \[
  1)\ y=\frac{2}{x},\qquad
  2)\ y=\frac{x}{2},\qquad
  3)\ y=-x^2+3,\qquad
  4)\ y=-3x+1.
  \]
  Типы:
  \begin{enumerate}
    \item убывающая прямая;
    \item гипербола с ветвями в I и III четвертях;
    \item парабола с ветвями вниз;
    \item возрастающая прямая.
  \end{enumerate}

  \item Для параболы
  \[
  y=x^2+6x+5
  \]
  определите знак \(a\), знак \(c\), а также найдите координаты вершины.

  \item Для гиперболы определите, какая из формул задаёт график, ветви которого расположены во II и IV четвертях и находятся ближе к осям, чем график \(y=\dfrac1x\):
  \[
  1)\ y=-\frac{3}{x},\qquad
  2)\ y=\frac{1}{3x},\qquad
  3)\ y=-\frac{1}{3x},\qquad
  4)\ y=\frac{3}{x}.
  \]

  \item На рисунке парабола имеет ветви вниз, пересекает ось \(Oy\) выше нуля, а её вершина находится левее оси \(Oy\). Определите знаки \(a\), \(b\), \(c\).

  \item Докажите, что если в треугольнике проведена медиана, то она делит треугольник на два равновеликих треугольника.

  \item В треугольнике \(ABC\) точка \(M\) --- середина \(AC\), точка \(K\) --- середина \(BM\). Найдите отношение
  \[
  S_{ABK}:S_{ABC}.
  \]

  \item В треугольнике \(ABC\) точка \(M\) --- середина \(AC\), точка \(K\) --- середина \(BM\). Через \(K\) проведена прямая, параллельная \(AB\), которая пересекает \(AC\) в точке \(J\), а \(BC\) --- в точке \(P\). Найдите
  \[
  S_{ABK}:S_{KPCM}.
  \]
\end{enumerate}

\bigskip
\textbf{Распределение по сложности:}
1--3 --- простые, 4--6 --- средние, 7--9 --- выше среднего, 10 --- сложная.

\newpage

% ---------------- ANSWERS ----------------
\section*{Ответы}

\begin{table}[h!]
\centering
\caption{Краткие ответы к тренировочному блоку}
\begin{tabularx}{\linewidth}{@{}cX@{}}
\toprule
\textbf{\#} & \textbf{Ответ} \\
\midrule
1 & \(\displaystyle k=-2,\ b=3,\) формула: \(\,y=-2x+3.\) \\
2 & \(\displaystyle a>0,\ b<0,\ c<0.\) \\
3 & \(\displaystyle x_{\text{в}}=1,\ y_{\text{в}}=3.\) \\
4 & \(1\!-\!\text{б},\ 2\!-\!\text{г},\ 3\!-\!\text{в},\ 4\!-\!\text{а}.\) \\
5 & \(\displaystyle a>0,\ c>0,\ x_{\text{в}}=-3,\ y_{\text{в}}=-4.\) \\
6 & \(\displaystyle y=-\frac{1}{3x}.\) \\
7 & \(\displaystyle a<0,\ b<0,\ c>0.\) \\
8 & Площади равны: \(\displaystyle S_{ABM}=S_{MBC}\), так как у треугольников общая высота к прямой \(AC\), а основания \(AM\) и \(MC\) равны. \\
9 & \(\displaystyle S_{ABK}:S_{ABC}=1:4.\) \\
10 & \(\displaystyle S_{ABK}:S_{KPCM}=4:7.\) \\
\bottomrule
\end{tabularx}
\end{table}

\vfill
\begin{center}
{\small\textcolor{darkaccent}{Памятка: сначала определяйте тип графика, затем --- знак коэффициентов, и только потом переходите к вычислениям.}}
\end{center}

\end{document}